% LuaLaTeXで2回コンパイルしてください。
\documentclass[a4paper,11pt]{ltjsarticle}
\usepackage[margin=24mm]{geometry}
\usepackage{amsmath,amssymb,booktabs}
\usepackage{hyperref}
\hypersetup{colorlinks=true,linkcolor=blue,urlcolor=blue}
\newcommand{\code}[1]{\texttt{\detokenize{#1}}}
\newcommand{\impl}[2]{\par\smallskip\noindent
 {\small 対応実装：\code{#1}、#2行目。}\par\smallskip}
\title{EpisodicDT-smallの学習方法と生成アルゴリズム\\
 \large 残差生成・終端SNRゼロ・v予測・CRPSによるモデル選択}
\author{}
\date{2026年9月8日版}
\begin{document}
\maketitle
\begin{abstract}
本資料は、過去のRSRPを条件として将来のRSRP分布を生成する現行モデルを説明する。
Transformer Encoderがhistoryを特徴ベクトルに変換し、条件付き拡散モデルが
history最終値からの変化量を生成する。重み更新にはv予測の平均二乗誤差とAdamWを用いる。
検証では固定した乱数で将来系列を生成し、CRPSで最良モデルを選択してEarly Stoppingを行う。
数式とソースコードの行番号を対応させて説明する。
\end{abstract}
\tableofcontents
\newpage

\section{目的と記号}
入力は観測済みのhistory、出力はその直後のfutureの複数候補である。
目的は条件付き分布$p_\theta(\mathbf y\mid\mathbf h)$の近似である。
本実装はRSRPの条件付き生成モデルであり、行動・報酬を用いた強化学習ではない。
\begin{center}
\begin{tabular}{lll}
\toprule 記号 & 意味 & 現在の設定\\ \midrule
$H$ & historyの観測数 & 40\\
$F$ & futureの観測数 & 10\\
$B$ & 学習バッチサイズ & 64\\
$D$ & Transformerの特徴次元 & 128\\
$C$ & contextの次元 & 32\\
$T$ & 拡散段階数 & 100\\
$S$ & 各検証historyの生成系列数 & 32\\ \bottomrule
\end{tabular}
\end{center}
futureの時刻位置$j=1,\ldots,F$と、拡散段階$t=1,\ldots,T$を区別する。
本資料の$t$はコードの\code{step}に1を加えた値である。
図とCSVの\code{forecast_step}は$j-1$である。
モデルはfutureの全$F$観測点をまとめて生成する。
生成した1観測点を入力として次の観測点を予測する自己回帰方式ではない。

\section{データ準備とエピソード}
対象はSUTDの次の4シナリオである。
\begin{itemize}
\item \code{Lvl4_AllRRUOn_Anomaly_label.csv}
\item \code{Lvl5_AllRRUOn_Anomaly_label.csv}
\item \code{Lvl6_1RRUOn_Anomaly_label.csv}
\item \code{Lvl6_AllRRUOn_Anomaly_label.csv}
\end{itemize}
\code{Time}と\code{RSRP}を抽出し、欠損行を除き、時刻順に整列する。
有効行数の前半を学習用、後半を最終評価用として保存する。出力列は\code{time}と\code{rsrp}である。
\impl{src/data/prepare_sutd_5g.py}{24--35}

学習用ファイルをさらに時系列順で80\%と20\%に分ける。
前方80\%で重みを更新し、後方20\%でモデルを選択する。
元データ全体では概ね40\%が重み更新、10\%が検証、50\%が最終評価となる。
丸めと境界ウィンドウ除外のため、エピソード数の比率はこれと異なる。

長さ$N$のトレースについて、開始位置$s$を0始まりとすると、
\[
 \mathbf h_s=(r_s,\ldots,r_{s+H-1}),\qquad
 \mathbf y_s=(r_{s+H},\ldots,r_{s+H+F-1}).
\]
開始位置を1点ずつずらし、$N-H-F+1$個の候補を作る。
検証境界$k=\lfloor0.8N\rfloor$に対し、
\[
 s+H+F\le k\Rightarrow\text{学習},\qquad s\ge k\Rightarrow\text{検証}
\]
とし、境界をまたぐ候補を除外する。各分割内ではエピソードが重複するが、
学習・検証間および異なるシナリオ間をまたぐエピソードは作らない。
\impl{src/data/dataset.py}{34--50}
\impl{src/train/train.py}{18--35}

約2 Hzの計測なので、history 40点は約20秒、future 10点は約5秒に相当する。
厳密な経過時間はCSV時刻で確認する必要がある。現在のdatasetはRSRP列だけを読み、
時刻間隔の補間や欠測区間の分断は行わない。

\subsection{標準化}
重み更新用の前方区間の観測値を連結した集合$\mathcal R_{\rm train}$から、
\[
 \mu=\frac1M\sum_{r\in\mathcal R_{\rm train}}r,\quad
 \sigma=\max\left(\sqrt{\frac1M\sum_{r\in\mathcal R_{\rm train}}(r-\mu)^2},10^{-6}\right)
\]
を計算する。$M$は観測数である。検証・最終評価の値やウィンドウの重複回数を使わない。
\[
 \tilde h_i=(h_i-\mu)/\sigma,\qquad \tilde y_j=(y_j-\mu)/\sigma.
\]
dataset構築時には一度全入力トレースから統計量を計算するが、
train.pyがデータ取得前に学習前方区間の統計量で上書きする。
検証と最終評価でも同じ$\mu,\sigma$を使う。
\impl{src/train/train.py}{60--72}

\section{Transformer Encoder}
入力形状は$[B,H,1]$である。各RSRPを線形変換し、位置埋め込みを加える。
\begin{align}
 \mathbf e_i&=W_{\rm in}\tilde h_i+\mathbf b_{\rm in}+\mathbf p_i,\\
 p_{i,2k}&=\sin(i/10000^{2k/D}),\\
 p_{i,2k+1}&=\cos(i/10000^{2k/D}).
\end{align}
位置埋め込みは固定式であり、学習パラメータではない。
\impl{src/models/episodic_diffusion.py}{12, 27--34}

各自己注意ヘッドは概念的に
\[
 \operatorname{Attention}(Q,K,V)=
 \operatorname{softmax}(QK^\top/\sqrt{d_k})V
\]
を計算する。4ヘッドなので$d_k=32$である。
自己注意・残差接続・LayerNorm・GELUを含むフィードフォワード処理を3層重ねる。
フィードフォワード内部次元は$4D=512$、Dropoutは0.05である。
因果マスクは指定していないため、history内の全位置を参照できる。
futureの正解値は入力しない。

最終historyトークンから、
\[
 \mathbf c=f_\phi(\tilde{\mathbf h})
 =\operatorname{SiLU}(W_c\mathbf e^{(L)}_{H-1}+\mathbf b_c)
 \in\mathbb R^C
\]
を得る。出力形状は$[B,32]$である。各成分に「傾き」等の意味を固定的には割り当てない。
\impl{src/models/episodic_diffusion.py}{13--22, 35--37}

\section{history最終値からの残差生成}
拡散対象のクリーンな系列を$\mathbf z_0$として、
\[
 z_{0,j}=\tilde y_j-\tilde h_{H-1}=\frac{y_j-h_{H-1}}{\sigma}
\]
を学習する。これは共通のhistory最終値との差であり、隣接future間の差分ではない。
history入力は絶対RSRPの標準化値のままである。
\impl{src/models/episodic_diffusion.py}{94--97}

生成残差から、
\[
 \hat{\tilde y}_j=\tilde h_{H-1}+\hat z_{0,j},\qquad
 \hat y_j=\sigma\hat{\tilde y}_j+\mu=h_{H-1}+\sigma\hat z_{0,j}
\]
で絶対RSRPを復元する。sample()が最終値を加算し、評価コードがdBmへ戻す。
CSVのRSRP列には絶対値が保存される。
\impl{src/models/episodic_diffusion.py}{135}
\impl{src/evaluation/evaluate.py}{16--26}

\section{終端SNRゼロの順方向拡散}
DDPMの基礎は文献\cite{ddpm}、終端SNRの背景は文献\cite{zero}を参照されたい。
以下は現行実装に対応する。
線形スケジュール$\beta_t^{\rm base}$を$10^{-4}$から$0.02$まで作り、
\[
 a_t=\sqrt{\prod_{u=1}^t(1-\beta_u^{\rm base})},\qquad
 s_t=(a_t-a_T)\frac{a_1}{a_1-a_T}
\]
と再スケーリングする。
\[
 \bar\alpha_0=1,\quad \bar\alpha_t=s_t^2,\quad
 \alpha_t=\bar\alpha_t/\bar\alpha_{t-1},\quad \beta_t=1-\alpha_t.
\]
各エピソードで$t$を一様に選び、
\[
 \epsilon\sim\mathcal N(\mathbf0,I),\qquad
 \mathbf z_t=\sqrt{\bar\alpha_t}\mathbf z_0+
              \sqrt{1-\bar\alpha_t}\boldsymbol\epsilon
\]
を直接計算する。順方向の全段階を順に実行する必要はない。
\impl{src/models/episodic_diffusion.py}{80--92, 99--102}

終端では$\bar\alpha_T=0$、$\beta_T=1$、$\mathbf z_T=\epsilon$となる。
信号対雑音比$\bar\alpha_T/(1-\bar\alpha_T)$はゼロであり、
学習時の終端分布と、純粋なノイズから始める生成時の初期分布が一致する。

\section{v予測と重み更新}
学習対象は
\[
 \mathbf v_t=\sqrt{\bar\alpha_t}\boldsymbol\epsilon-
             \sqrt{1-\bar\alpha_t}\mathbf z_0
\]
である。velocityという名称は拡散のパラメータ化であり、
物理的な速度やRSRPの時間微分ではない。
\[
 \hat{\mathbf v}_t=g_\psi(\mathbf z_t,(t-1)/(T-1),\mathbf c).
\]
クラス名はNoisePredictorだが、出力はvである。拡散時刻を128次元へ埋め込み、
ノイズ付き残差10次元とcontext 32次元に連結する。
予測MLPは170次元入力、128次元隠れ層、10次元出力を持つ。
\impl{src/models/episodic_diffusion.py}{40--61, 103--107}

学習損失は
\[
 \mathcal L_v(\phi,\psi)=
 \mathbb E_{(\mathbf h,\mathbf y),t,\epsilon}
 \left[\frac1F\|\hat{\mathbf v}_t-\mathbf v_t\|_2^2\right].
\]
各バッチでは$BF$要素のMSEを計算する。エンコーダの別損失はなく、
\[
 \nabla_\phi\mathcal L_v
 =\left(\frac{\partial\mathbf c}{\partial\phi}\right)^\top
  \left(\frac{\partial\hat{\mathbf v}_t}{\partial\mathbf c}\right)^\top
  \nabla_{\hat{\mathbf v}_t}\mathcal L_v
\]
によって入力射影、自己注意、フィードフォワード、LayerNorm、context射影も更新される。
\impl{src/models/episodic_diffusion.py}{107--108}

逆伝播後、勾配の全体ノルムを1.0でクリッピングし、AdamWを適用する。
全パラメータを$\theta=(\phi,\psi)$、クリッピング後の勾配を$g_k$とすると、
\begin{align}
 m_k&=\beta_1m_{k-1}+(1-\beta_1)g_k,\\
 u_k&=\beta_2u_{k-1}+(1-\beta_2)g_k^2,\\
 \hat m_k&=m_k/(1-\beta_1^k),\quad \hat u_k=u_k/(1-\beta_2^k),\\
 \theta_{k+1}&=(1-\eta\lambda)\theta_k-
             \eta\hat m_k/(\sqrt{\hat u_k}+\varepsilon_{\rm opt}).
\end{align}
演算は要素ごとである。学習率$\eta=10^{-4}$以外のAdamW引数は実行環境の既定値を使う。
train\_v\_mseはバッチ損失をエピソード数で重み付けして平均した値で、dBm誤差ではない。
\impl{src/train/train.py}{78, 91--99}

\section{生成：vからの逆拡散}
historyからcontextを一度計算し、各historyについて$S$本の初期ノイズ
$\mathbf z_T^{(s)}\sim\mathcal N(\mathbf0,I)$を作る。
$t=T,T-1,\ldots,1$の各段階で、
\[
 \hat{\mathbf z}_0=\sqrt{\bar\alpha_t}\mathbf z_t-
                  \sqrt{1-\bar\alpha_t}\hat{\mathbf v}_t
\]
によりクリーンな残差を推定する。これは$\alpha_T=0$でもゼロ除算を起こさない。
DDPM事後分布の係数を用いて、
\begin{align}
 A_t&=\frac{\beta_t\sqrt{\bar\alpha_{t-1}}}{1-\bar\alpha_t},\\
 B_t&=\frac{(1-\bar\alpha_{t-1})\sqrt{\alpha_t}}{1-\bar\alpha_t},\\
 \tilde\beta_t&=\frac{\beta_t(1-\bar\alpha_{t-1})}{1-\bar\alpha_t},\\
 \mathbf z_{t-1}&=A_t\hat{\mathbf z}_0+B_t\mathbf z_t+
 \mathbf1[t>1]\sqrt{\tilde\beta_t}\boldsymbol\xi_t,\quad
 \boldsymbol\xi_t\sim\mathcal N(\mathbf0,I)
\end{align}
と更新する。係数は固定bufferとして保存する。最終更新$t=1$で追加ノイズは入れない。
\impl{src/models/episodic_diffusion.py}{84--92, 127--134}
最後にhistory最終値を加算する。sample()の出力は$[B,S,F,1]$の絶対RSRP標準化値である。
生成時はeval()でDropoutを無効化し、no\_grad()で勾配計算を省く。
実測futureは生成に入力せず、出力後の比較にのみ使う。
\impl{src/models/episodic_diffusion.py}{110--135}

\section{検証とモデル選択}
\subsection{CRPS}
検証エピソード$n$、future位置$j$の実測を$y_{n,j}$、生成を$\hat y_{n,j}^{(s)}$とする。
経験分布のCRPSは
\[
 \operatorname{CRPS}_{n,j}=
 \frac1S\sum_s|\hat y^{(s)}_{n,j}-y_{n,j}|-
 \frac1{2S^2}\sum_s\sum_r|\hat y^{(s)}_{n,j}-\hat y^{(r)}_{n,j}|.
\]
第1項は実測との距離、第2項は分布の広がりを考慮する。
全検証エピソード・全future位置で平均した値をモデル選択に使う。小さいほどよく、単位はdBmである。
CRPSでは逆伝播しない。重み更新はv-MSE、モデル選択は生成後のCRPSである。

ソートした生成値$\hat y_{(s)}$を使うと、上式第2項は
\[
 S^{-2}\sum_{s=1}^S(2s-S-1)\hat y_{(s)}
\]
となり、$S\times S$の距離配列を保持せずに計算できる。
\impl{src/train/validation.py}{6--19, 32--36}

\subsection{ステップ別指標}
生成平均を$\bar y_{n,j}=S^{-1}\sum_s\hat y^{(s)}_{n,j}$とすると、
\begin{align}
 \operatorname{MAE}_j&=N_{\rm val}^{-1}\sum_n|\bar y_{n,j}-y_{n,j}|,\\
 \operatorname{Bias}_j&=N_{\rm val}^{-1}\sum_n(\bar y_{n,j}-y_{n,j}).
\end{align}
各$(n,j)$の生成サンプルの5\%・95\%分位点を$q^{.05}_{n,j},q^{.95}_{n,j}$とすると、
\[
 \operatorname{Coverage}_{90,j}=N_{\rm val}^{-1}\sum_n
 \mathbf1[q^{.05}_{n,j}\le y_{n,j}\le q^{.95}_{n,j}].
\]
区間幅、実測標準偏差、生成標準偏差も記録する。
生成標準偏差はエピソードとサンプルを両方まとめた分布の値であり、
単一エピソード内の不確実性だけを表すものではない。
経験的90\%帯は被覆率90\%を保証しない。
\impl{src/train/validation.py}{20--36}

\subsection{固定ノイズとEarly Stopping}
検証損失はシード12345、生成は12346を使い、毎エポック同じ乱数条件で比較する。
fork\_rngにより検証後は学習側の乱数状態を復元する。
同一環境・同一バッチ構成での再現を意図し、異なるGPU等でのビット単位一致を保証するものではない。
\impl{src/train/validation.py}{40--67}

保存と停止は別の条件で判定する。
\begin{enumerate}
\item CRPSの厳密な最小値を更新すれば、その場で重みを保存する。
\item 停止用基準値より0.001 dBmを超えて改善すれば、基準値を更新して非改善カウンタを0にする。
それ以外はカウンタを1増やす。
\item 30エポック以上実行し、カウンタが15以上なら停止する。上限は150エポックである。
\end{enumerate}
0.001 dBm未満の改善も最良モデルの保存対象だが、停止カウンタを毎回リセットするわけではない。
選ばれるのは検証基準での最良エポックであり、未知データでの絶対的最適値を保証しない。
\impl{src/train/train.py}{100--138}

\section{アルゴリズムと出力}
\subsection{学習の手順}
\begin{enumerate}
\item 学習・検証を時系列分割し、学習区間でのみ$\mu,\sigma$を計算する。
\item モデルとAdamWを初期化する。
\item 各エポックで学習用エピソードをシャッフルし、ミニバッチを取得する。
\item historyを符号化し、futureをhistory最終値からの残差に変換する。
\item ランダムな$t,\epsilon$から$\mathbf z_t,\mathbf v_t$を作る。
\item v-MSEを逆伝播し、勾配クリッピング後にAdamWで更新する。
\item 固定乱数で全検証historyから生成し、CRPS等を計算する。
\item 最良重みとログを保存し、停止条件を判定する。
\end{enumerate}

\subsection{最終評価}
保存形式を確認し、チェックポイントから構造・重み・標準化統計量を復元する。
各評価historyから逆拡散し、全生成futureと実測futureをCSVに保存する。
集約箱ひげ図、ステップ別箱ひげ図、small multiplesも出力する。
\impl{src/evaluation/evaluate.py}{137--159, 176--211}

最終評価の全体MAEは生成平均によるMAEであり、small multiplesのタイトルの
生成中央値によるMAEとは異なる。コンソールの危険確率は現状では最後のエピソードだけの値である。
CRPSとステップ別数値ログは学習時の検証で計算する。
\impl{src/evaluation/evaluate.py}{161--172}

\section{設定と実行}
現在の設定値はhistory 40、future 10、特徴次元128、context次元32、
4ヘッド・3層・Dropout 0.05、拡散100段階、学習率0.0001、バッチ64である。
検証生成数32、検証割合0.2、最低30・最大150エポック、
patience 15、min\_delta 0.001 dBmである。
\begin{verbatim}
python -m src.data.prepare_sutd_5g
python -m src.train.train --config configs/config.yaml
python -m src.evaluation.evaluate --config configs/config.yaml --samples 32
\end{verbatim}
モデルは\code{results/checkpoints/residual_v1/episodicdt.pt}に、
ログは同じディレクトリの\code{training_metrics.json}に保存する。
モデル形式は\code{zero_snr_v_residual_v1}で、旧epsilon予測の重みはそのまま使用できない。
評価時の構造とhistory/future長は保存設定から、評価CSVパスは実行時設定から取得する。

\section{資料の範囲と検証}
行番号は2026年9月8日に確認したソースに対応する。コード変更後は行番号も変わる。
残差復元、終端の有限値、エンコーダ勾配、CRPS式の一致、乱数復元、
学習・停止・評価出力のテストを用意している。
\begin{verbatim}
python -m unittest discover -s tests -v
\end{verbatim}
これは実装の動作確認であり、SUTD実データの精度向上を示す実験結果ではない。
ステップ別CRPSや箱ひげ図は各時刻の分布比較であり、
時系列の変化パターンや時刻間相関の一致を単独では保証しない。

\begin{thebibliography}{9}
\bibitem{ddpm} J. Ho, A. Jain, P. Abbeel,
\emph{Denoising Diffusion Probabilistic Models}, 2020.
\url{https://arxiv.org/abs/2006.11239}
\bibitem{zero} S. Lin, B. Liu, J. Li, X. Yang,
\emph{Common Diffusion Noise Schedules and Sample Steps are Flawed}, 2023.
\url{https://arxiv.org/abs/2305.08891}
\end{thebibliography}
\end{document}
